% Revize 2026-10-28: obsah, matematika a návaznost; podrobnosti a časy
% viz planovani/REVIZE_PRIPRAV_1_ROCNIK.md.
% Aktuální porovnání se skutečným obsahem místních PDF Realisticky.cz:
% externi_zdroje/realisticky/01_Zakladni_poznatky/02_Ciselne_obory/010203_Racionalni_cisla_I/010203_Lekce.pdf
% externi_zdroje/realisticky/01_Zakladni_poznatky/02_Ciselne_obory/010204_Racionalni_cisla_II/010204_Lekce.pdf
% Rozpis a plán: 1_rocnik/tematicke_plany/Rozpis_vyuky_MAT_1_rocnik.tex,
% hodina 6; SPS.pdf, oddíl 1.2. Návaznost: hodiny 05 a 07, konvence N jako v 12.
% Učebnice: ucebnice/01_Zakladni_poznatky/Ciselne_obory/Číselné obory.tex,
% Racionální čísla, Zápisy racionálních čísel, Porovnávání racionálních čísel.
% Metodická inspirace: místní Realisticky.cz, pod externi_zdroje/realisticky/:
% 01_Zakladni_poznatky/01_Opakovani/010108_Zlomky_I/010108_Lekce.pdf (Př. 3–5);
% 01_Zakladni_poznatky/02_Ciselne_obory/010203_Racionalni_cisla_I/010203_Lekce.pdf
% (dělení nulou, základní tvar a různé zápisy);
% 01_Zakladni_poznatky/02_Ciselne_obory/010204_Racionalni_cisla_II/010204_Lekce.pdf
% (úvod a Př. 3, účelné krácení).
% Prostudovány uvedené části, nikoli převzat soubor úloh. Všechny úlohy jsou vlastní.
% Chybná položka -126/-198 v řetězci ekvivalentních zlomků v naší učebnici
% nebyla použita: její hodnota je 7/11, nikoli 2/3. Zdroj ponechán beze změny.
% Cílená revize: vlastní výklad a úlohy V1–V3, R2 k rozšíření Z -> Q,
% opačnému/převrácenému prvku a distributivitě podle zadání učitele.
% Podklady: v uvedené kapitole učebnice Vlastnosti racionálních čísel
% a Převrácené číslo; místní Realisticky 010203, Př. 1–2 a převrácené číslo,
% 010204, Př. 1–3. Důsledně rozlišena oboustranná neutralita a dělení nulou.
\documentclass[11pt,a4paper]{article}
% Spolecny styl pro ucitelske pripravy MAT (Work/MAT)
% Verze 1.1, 2026-10-23
\NeedsTeXFormat{LaTeX2e}
\ProvidesPackage{priprava}[2026/10/23 v1.1 Sablona ucitelskych priprav MAT]

% Zakladni typografie a cestina
\RequirePackage[T1]{fontenc}
\RequirePackage[utf8]{inputenc}
\RequirePackage[czech]{babel}
\RequirePackage{lmodern}
\RequirePackage{microtype}
\RequirePackage{amsmath,amssymb}
\RequirePackage{xcolor}
\RequirePackage{enumitem}
\RequirePackage[most]{tcolorbox}
\RequirePackage{tabularx,booktabs,array}
\RequirePackage{fancyhdr}
\RequirePackage{needspace}
\RequirePackage{geometry}
\RequirePackage[hidelinks]{hyperref}

\geometry{a4paper,top=17mm,bottom=18mm,left=23mm,right=23mm}
\definecolor{notegray}{gray}{0.39}
\setlength{\parindent}{0pt}
\setlength{\parskip}{2.5pt plus 1pt minus 0.5pt}
\setlist[itemize]{topsep=2pt,itemsep=1pt,parsep=0pt,leftmargin=1.45em}
\setlist[enumerate]{topsep=2pt,itemsep=2pt,parsep=0pt,leftmargin=1.65em}
\setcounter{secnumdepth}{0}
\renewcommand{\arraystretch}{1.12}

% Zahlavi: vlevo rocnik | globalni tema, vpravo cislo: nazev hodiny
% Zapati zustava zcela prazdne.
\pagestyle{fancy}
\fancyhf{}
\renewcommand{\headrulewidth}{0.35pt}
\renewcommand{\footrulewidth}{0pt}
\setlength{\headheight}{14pt}
\newcommand{\nastavPripravu}[4]{%
  \fancyhead[L]{\small #1\enspace\textbar\enspace #2}%
  \fancyhead[R]{\small #3: #4}%
}

% Minimalisticke cernobile boxy. Kazda definice zvlast.
\tcbset{commonlesson/.style={enhanced,breakable=false,colback=white,colframe=black,
  colbacktitle=white,coltitle=black,boxrule=0.45pt,
  arc=0.8mm,outer arc=0.8mm,left=3mm,right=3mm,top=1.6mm,bottom=1.6mm,
  toptitle=1mm,bottomtitle=0.7mm,fonttitle=\bfseries,
  attach boxed title to top left={xshift=2.5mm,yshift=-1.15mm},
  boxed title style={colback=white,colframe=white,boxrule=0pt,left=1mm,right=1mm}}}
\newtcolorbox{zadaniBox}[1][]{commonlesson,title={Zadání},#1}
\newtcolorbox{prikladBox}[1][]{commonlesson,title={Příklad},#1}
\newtcolorbox{definiceBox}[1][]{commonlesson,title={Definice},#1}

% Spolecny box pro zadani a jeho reseni, bez mezery mezi castmi.
% Pouziti:
%   \begin{prikladReseniBox}
%     Zadani ulohy.
%     \tcblower
%     \begin{align*} ... \end{align*}
%   \end{prikladReseniBox}
% Spodni cast ma jemny sedy podklad a delici linku; lze v ni sazet
% vice radku matematiky. Obe casti maji jen jeden spolecny ramecek.
\tcbset{solutionlesson/.style={
  commonlesson,
  enhanced jigsaw,
  breakable,
  colbacklower=black!3,
  segmentation style={solid,draw=black!25,line width=0.35pt},
  fontlower=\small\color{notegray},
  before lower={\noindent\textbf{Řešení:}\enspace},
  leftlower=3mm,rightlower=3mm
}}
\newtcolorbox{prikladReseniBox}[1][]{solutionlesson,title={Příklad},#1}
\newtcolorbox{zadaniReseniBox}[1][]{solutionlesson,title={Zadání},#1}

% Pomucky pro autora pripravy: sedive poznamky; \reseni pro starsi pripravy
\newcommand{\poznamka}[1]{\par\begingroup\small\color{notegray}\itshape
  \textnormal{\textbf{Poznámka pro učitele.}}\enspace #1\par\endgroup}
\newcommand{\reseni}[1]{\par\begingroup\small\color{notegray}
  \textbf{Očekávané řešení:}\enspace #1\par\endgroup}
\newcommand{\blok}[1]{\needspace{5\baselineskip}\subsection*{#1}}
\newcommand{\vyrok}[1]{\uv{#1}}

\nastavPripravu{1. ročník}{Číselné obory}{06}{Racionální čísla a zlomky}
\begin{document}

\poznamka{Rozšířená příprava obsahuje více materiálu než na jedinou
hodinu. Pro 45 minut orientačně vyber: význam zlomku a rozšíření oboru
6 min, ekvivalentní zápisy 8 min, opačné a převrácené číslo 9 min,
distributivita 7 min, porovnávání a vybrané samostatné úlohy 13 min,
shrnutí 2 min. Další procvičení je zásoba; původní D1 je asi na 15–20 min.
Rezervu zadávej výběrově. Přednost má vysvětlení kroků před dlouhým počítáním.}

\blok{Jedno číslo může mít více zápisů}
Tři stejné tyče rozdělíme rovným dílem mezi pět lidí. Každý dostane
$\frac35$ tyče. Rozdělením každé pětiny ještě napůl dostane každý
šest desetin téže tyče. Zápisy $\frac35$ a $\frac6{10}$ označují stejné číslo.

\begin{definiceBox}
\textbf{Racionální číslo} je číslo zapsatelné jako $\frac pq$, kde
$p,q\in\mathbb Z$ a $q\ne0$. Obor racionálních čísel značíme $\mathbb Q$.
\end{definiceBox}

Ve zlomku $\frac pq$ je $p$ čitatel, $q$ jmenovatel a zlomková čára
znamená dělení. Každé celé číslo je racionální: například $-4=\frac{-4}{1}$.
Zlomek s nulovým čitatelem a nenulovým jmenovatelem je nula,
například $\frac07=0$. Zápisy $\frac50$ ani $\frac00$ číslo nevyjadřují.
Pro přehlednost budeme záporné znaménko psát před zlomek a jmenovatel kladný.

\begin{definiceBox}
\textbf{Rozšiřování zlomku} znamená vynásobení čitatele i jmenovatele
týmž nenulovým celým číslem $k$:
$\frac pq=\frac{kp}{kq}$. Hodnota se nemění, protože násobíme číslem $\frac kk=1$.
\end{definiceBox}

\begin{definiceBox}
\textbf{Krácení zlomku} znamená vydělení čitatele i jmenovatele
týmž jejich nenulovým společným celočíselným dělitelem. Hodnota zlomku se nemění.
\end{definiceBox}

\begin{definiceBox}
Zlomek s kladným jmenovatelem je v \textbf{základním tvaru}, jestliže
čitatel a jmenovatel nemají společného kladného dělitele většího než $1$.
\end{definiceBox}

\begin{prikladReseniBox}[title={P1: zjednodušení a kontrola},breakable=false]
Zkrať do základního tvaru $\frac{-84}{126}$ a $\frac{-45}{-60}$.
První výsledek rozšiř na jmenovatele $30$.
\tcblower
$\frac{-84}{126}=-\frac{2\cdot42}{3\cdot42}=-\frac23$;
$\frac{-45}{-60}=\frac{45}{60}=\frac34$.
Rozšíření: $-\frac23=-\frac{20}{30}$. Zpětně $(-2)\cdot126=(-84)\cdot3=-252$.
U druhého zlomku dvě záporná znaménka dávají kladný podíl.
\end{prikladReseniBox}

\poznamka{U zlomku $0/q$ vede základní tvar k $0/1$.
Krátíme společné \emph{činitele}, ne podobně vypadající číslice ani části součtu.
Čitatel a jmenovatel nesčítáme, chceme-li zlomek rozšířit.}

\newpage
\blok{Proč při dělení nevystačíme s celými čísly}
V $\mathbb Z$ lze sčítat, odčítat a násobit bez opuštění oboru.
U dělení to nestačí ani po vyloučení nuly z dělitelů: některý podíl
celý je, jiný není. Racionální čísla doplní všechny podíly celých čísel
s nenulovým dělitelem.

\begin{prikladReseniBox}[title={V1: některé dělení vyjde v Z, jiné až v Q},breakable=false]
Vypočti $12:3$, $5:2$, $(-7):3$ a $0:4$.
Které výsledky jsou celé? Výsledky ověř násobením dělitelem.
\tcblower
$12:3=4\in\mathbb Z$ a $0:4=0\in\mathbb Z$.
Naproti tomu $5:2=\frac52=2{,}5\notin\mathbb Z$ a
$(-7):3=-\frac73\notin\mathbb Z$; oba podíly jsou racionální.
Zkoušky: $3\cdot4=12$, $2\cdot\frac52=5$,
$3\cdot(-\frac73)=-7$ a $4\cdot0=0$.
Věta \uv{V celých číslech dělení nikdy nevyjde} je tedy stejně chybná
jako tvrzení, že vyjde vždy.
\end{prikladReseniBox}

\blok{Opačné číslo sčítáme, převráceným násobíme}
\begin{definiceBox}
Připomenutí: \textbf{opačné číslo} k $a\in\mathbb Q$ je $-a$;
platí $a+(-a)=0$. Jde o opačný prvek vzhledem ke sčítání.
\end{definiceBox}

\begin{definiceBox}
\textbf{Převrácené číslo} k nenulovému $a\in\mathbb Q$ je číslo $\frac1a$,
jehož součin s $a$ je $1$. Jde o inverzní prvek vzhledem k násobení:
$a\cdot\frac1a=1$.
\end{definiceBox}

Proč existuje pro každé nenulové racionální číslo? Zapíšeme
$a=\frac pq$ s celými $p,q$, kde $q\ne0$. Z předpokladu $a\ne0$ plyne
také $p\ne0$. Číslo $\frac qp$ je proto opět racionální a
\[
\frac pq\cdot\frac qp=\frac{pq}{qp}=1.
\]
Nula má opačné číslo $0$, ale převrácené číslo nemá:
pro žádné $b$ neplatí $0\cdot b=1$.

\begin{prikladReseniBox}[title={V2: dva různé návraty k neutrálnímu prvku},breakable=false]
Doplň opačné a převrácené číslo k $-\frac34$, $\frac25$, $-1$ a $0$.
U prvního čísla proveď zkoušku součtem i součinem.
\tcblower
\begin{center}
\begin{tabular}{ccc}
\toprule
Číslo $a$ & Opačné $-a$ & Převrácené $\frac1a$, pokud existuje\\
\midrule
$-\frac34$ & $\frac34$ & $-\frac43$\\
$\frac25$ & $-\frac25$ & $\frac52$\\
$-1$ & $1$ & $-1$\\
$0$ & $0$ & neexistuje\\
\bottomrule
\end{tabular}
\end{center}
Zkoušky: $-\frac34+\frac34=0$ a
$(-\frac34)\cdot(-\frac43)=1$. Opačné číslo u nenulového čísla
změní znaménko; převrácené číslo je stejného znaménka jako původní číslo.
\end{prikladReseniBox}

V $\mathbb Z$ nemá například $2$ převrácený prvek pro násobení,
protože $\frac12$ není celé číslo. V $\mathbb Q$ už jej má.
Pro $b\ne0$ tak dělení nahradíme násobením převráceným číslem:
například $\frac34:\frac25=\frac34\cdot\frac52=\frac{15}{8}$.
Zkouška $\frac{15}{8}\cdot\frac25=\frac34$ vrací původní dělenec.

\poznamka{U každého \uv{inverzního} prvku chtěj doplnit operaci
a její neutrální prvek: pro součet jde o $0$, pro součin o $1$.
Dělení samo není komutativní: $6:3=2$, ale $3:6=\frac12$.
Jednička není oboustranně neutrální pro dělení: $6:1=6$, ale $1:6\ne6$.
To je jiný výrok než neutralita jedničky pro násobení.}

\blok{Distributivita propojuje násobení se sčítáním}
\begin{definiceBox}
\textbf{Distributivita násobení vzhledem ke sčítání} znamená, že pro
všechna $a,b,c\in\mathbb Q$ platí
\[
a\cdot(b+c)=a\cdot b+a\cdot c.
\]
Stejně platí $(b+c)\cdot a=b\cdot a+c\cdot a$.
\end{definiceBox}

Pravidlo platí už v $\mathbb N$ a $\mathbb Z$, nyní také pro zlomky.
Šest stejných dvojic dílů můžeme spočítat jako šestkrát součet,
nebo zvlášť šestkrát každý díl a výsledky sečíst. Podstatné je,
že vně závorky je násobení, uvnitř sčítání.

\begin{prikladReseniBox}[title={V3: dvě cesty k témuž výsledku},breakable=false]
Vypočti dvojím způsobem $6\cdot(\frac56-\frac13)$ a
$-\frac32\cdot(\frac23+\frac49)$. Potom posuď chybný přepis
$6:(2+1)=6:2+6:1$.
\tcblower
Nejprve závorky, potom roznásobení:
\begin{align*}
6\cdot\left(\frac56-\frac13\right)&=6\cdot\frac12=3,\\
6\cdot\frac56-6\cdot\frac13&=5-2=3,\\
-\frac32\cdot\left(\frac23+\frac49\right)&=-\frac32\cdot\frac{10}{9}=-\frac53,\\
-\frac32\cdot\frac23-\frac32\cdot\frac49&=-1-\frac23=-\frac53.
\end{align*}
Odčítání v první závorce je sčítání opačného čísla, takže pravidlo platí
i v tomto zápisu. U dělení součtem přepis neplatí: vlevo je $6:3=2$,
vpravo $3+6=9$. Dělení každým sčítancem je jiná operace než dělení součtem.
\end{prikladReseniBox}

\poznamka{Příklad s roznásobením nevyužívá žádný nový vzorec pro
algebraické výrazy. Zachovej konkrétní zlomky, nech žáky porovnat oba
postupy. Samotná shoda příkladů ilustruje obecné pravidlo, nedokazuje je.}

\blok{Proč pravidla fungují i při úpravách zlomků}
Společný jmenovatel znamená stejně velké díly. Například třetinu
a polovinu stejného celku přepíšeme na šestiny:
\[
\frac13+\frac12=\frac26+\frac36=\frac{2+3}{6}=\frac56.
\]
Poslední krok je distributivita v opačném směru:
$2\cdot\frac16+3\cdot\frac16=(2+3)\cdot\frac16$.
Krácení zase používá násobení převrácenou hodnotou společného
nenulového činitele, nikoli odstraňování sčítanců.
Proto lze například
$\frac{3\cdot2+3\cdot5}{3\cdot4}
=\frac{3(2+5)}{3\cdot4}=\frac74$;
nejprve vytkneme společný činitel z \emph{celého} čitatele.

Pro celá $p,r,q,s$ s $q\ne0$, $s\ne0$ dostáváme obecně
\begin{align*}
\frac pq\mathbin{\pm}\frac rs&=\frac{ps\pm rq}{qs},
&\frac pq\cdot\frac rs&=\frac{pr}{qs},\\
\frac pq:\frac rs&=\frac pq\cdot\frac sr=\frac{ps}{qr}
&&\text{navíc }r\ne0.
\end{align*}
Čitatelé i jmenovatelé výsledků jsou celí a uvedené jmenovatele
jsou nenulové. Tím jsme zdůvodnili uzavřenost $\mathbb Q$ pro součet,
rozdíl, součin a dělení nenulovým číslem, ne pouze ověřili několik příkladů.
Převrácené číslo je také jediné: kdyby $ab=1$, pak násobením
rovnosti číslem $1/a$ dostaneme $b=1/a$.

\poznamka{Navazuje na Realisticky 010204, Př. 1–3: společné díly,
odvození obecných pravidel a krácení před násobením.
Ve třídě stačí odvodit součet a podíl; další zápisy slouží jako přehled.
Tento rozbor nahrazuje část mechanických výpočtů V3 nebo S1,
nepřidává se povinně k celé sadě.}

\newpage
\blok{Porovnávání bez odhadování}
Při stejném kladném jmenovateli rozhoduje čitatel. Při různých kladných
jmenovatelích zlomky rozšíříme na společný jmenovatel. Pro $b,d>0$
tedy platí $\frac ab<\frac cd$ právě tehdy, když $ad<bc$.
Podmínka kladných jmenovatelů brání chybě se znaménkem nerovnosti.

\begin{prikladReseniBox}[title={P2: záporná čísla a společný jmenovatel},breakable=false]
Porovnej $-\frac56$ a $-\frac78$. Potom seřaď vzestupně
$-\frac56$, $-\frac78$, $-\frac34$, $0$.
\tcblower
$-\frac56=-\frac{20}{24}$, $-\frac78=-\frac{21}{24}$,
$-\frac34=-\frac{18}{24}$. Proto
\[
-\frac78<-\frac56<-\frac34<0.
\]
Číslo $-21$ je menší než $-20$; záporné znaménko nesmíme při porovnávání ignorovat.
\end{prikladReseniBox}

\begin{prikladReseniBox}[title={P3: desetinný zápis a číslo mezi zlomky},breakable=false]
Zapiš $0{,}375$ jako zlomek v základním tvaru.
Najdi racionální číslo ostře mezi $\frac23$ a $\frac34$.
\tcblower
$0{,}375=\frac{375}{1000}=\frac38$.
Dále $\frac23=\frac{16}{24}$ a $\frac34=\frac{18}{24}$,
takže vyhovuje $\frac{17}{24}$. Nejde o jediné možné číslo.
\end{prikladReseniBox}

\blok{Samostatně: zápis, velikost, časté chyby}
\begin{zadaniReseniBox}[title={S1: základní sada},breakable=false]
\begin{enumerate}[label=\alph*)]
\item Zkrať $\frac{54}{90}$, $\frac{72}{-120}$, $\frac{-63}{-84}$.
\item Doplň $\frac{-3}{5}=\frac{\square}{35}$,
$\frac7{12}=\frac{42}{\square}$ a $\frac07=\frac{\square}{1}$.
\item Porovnej $\frac7{10}$ a $\frac{11}{15}$;
$-\frac49$ a $-\frac5{12}$; $\frac{-5}{-8}$ a $\frac35$.
\item Zapiš $-1{,}25$ jako zlomek a $\frac7{20}$ desetinným číslem.
\item Najdi jedno racionální číslo mezi $\frac35$ a $\frac23$.
\end{enumerate}
\tcblower
a) $\frac35$, $-\frac35$, $\frac34$.
b) $-21$, $72$, $0$.
c) $\frac{21}{30}<\frac{22}{30}$;
$-\frac{16}{36}<-\frac{15}{36}$;
$\frac58=\frac{25}{40}>\frac{24}{40}=\frac35$.
d) $-\frac{125}{100}=-\frac54$; $0{,}35$.
e) Například $\frac{19}{30}$, protože $\frac35=\frac{18}{30}$ a $\frac23=\frac{20}{30}$.
\end{zadaniReseniBox}

\begin{zadaniReseniBox}[title={S2: oprav a vysvětli},breakable=false]
Rozhodni, zda platí: a) $\frac23=\frac{2+3}{3+3}$;
b) $\frac{2+6}{2+10}=\frac6{10}$;
c) $\frac{6\cdot5}{6\cdot7}=\frac57$.
\tcblower
a) Ne, $\frac56\ne\frac23=\frac46$: přičtení stejného čísla není rozšíření.
b) Ne, vlevo je $\frac8{12}=\frac23$, vpravo $\frac35$.
Nelze krátit sčítance. c) Ano, krátíme společného činitele $6$.
\end{zadaniReseniBox}

\poznamka{Při rozboru chtěj vedle výsledku vysvětlení: která společná
jednotka se použila a proč znaménko vyšlo právě takto. Závěrečná otázka:
\uv{Je každý zlomek jiné číslo?} Hodina 07 rozšíří obor o čísla,
která jako podíl celých čísel vyjádřit nejdou.}

\newpage
\blok{Rezerva}
\begin{zadaniReseniBox}[title={R1: volitelně, 10–15 minut},breakable=false]
\begin{enumerate}[label=\alph*)]
\item Porovnej bez desetinného dělení $\frac{101}{102}$ a $\frac{100}{101}$.
\item Urči všechna přirozená $n$, pro která $\frac23<\frac n{12}<\frac56$.
\item Tři čtvrtiny objemu nádrže tvoří $18$ litrů. Kolik litrů
má celá nádrž a kolik litrů představuje jedna šestina jejího objemu?
\item Vysvětli, proč \uv{větší jmenovatel znamená menší zlomek} bez dalších
podmínek nestačí. Uveď dva kladné zlomky jako protipříklad.
\end{enumerate}
\tcblower
a) První je $1-\frac1{102}$, druhý $1-\frac1{101}$.
Protože $\frac1{102}<\frac1{101}$, je první zlomek větší.
Také $101\cdot101=10\,201>10\,200=100\cdot102$.
b) $\frac8{12}<\frac n{12}<\frac{10}{12}$, tedy pouze $n=9$.
c) Jedna čtvrtina je $6$ litrů, celek $24$ litrů, šestina $4$ litry.
d) Například $\frac45>\frac12$, ačkoli $5>2$.
Pravidlo platí při stejném kladném čitateli a kladných jmenovatelích.
\end{zadaniReseniBox}

\begin{zadaniReseniBox}[title={R2: volitelně k novým vlastnostem},breakable=false]
\begin{enumerate}[label=\alph*)]
\item Urči opačné a převrácené číslo k $-\frac58$ a proveď obě zkoušky.
\item Vypočti pomocí roznásobení $-\frac23\cdot(\frac94-\frac32)$.
Ověř výpočet nejprve sečtením v závorce.
\item Najdi všechna celá nenulová čísla, která mají také celočíselnou
převrácenou hodnotu. Zdůvodni úplnost výčtu.
\end{enumerate}
\tcblower
a) Opačné číslo je $\frac{5}{8}$ a převrácené číslo je $-\frac{8}{5}$.
Zkoušky:
\[
-\frac{5}{8}+\frac{5}{8}=0,\qquad
\left(-\frac{5}{8}\right)\cdot\left(-\frac{8}{5}\right)=1.
\]
b) $-\frac32+1=-\frac12$; kontrola $-\frac23\cdot\frac34=-\frac12$.
c) Právě $1$ a $-1$. Součin dvou celých čísel může být $1$ pouze jako
$1\cdot1$ nebo $(-1)\cdot(-1)$; celočíselní dělitelé jedničky jsou jen $1$ a $-1$.
\end{zadaniReseniBox}

\poznamka{R2 je volitelná rezerva nebo náhrada části následujícího D1,
nikoli povinný přídavek k celé domácí práci.}

\blok{Domácí procvičování}
\begin{zadaniReseniBox}[title={D1: přibližně 15–20 minut},breakable=false]
\begin{enumerate}[label=\arabic*)]
\item Zkrať $\frac{96}{144}$, $\frac{-75}{105}$, $\frac{-56}{-98}$.
\item Rozšiř $-\frac7{12}$ na jmenovatele $60$ a $\frac9{14}$ na $42$.
\item Porovnej $\frac58$ a $\frac{11}{18}$; $-\frac7{10}$ a $-\frac23$.
\item Seřaď vzestupně $\frac34$, $-\frac12$, $0{,}7$, $-\frac58$.
\item Zapiš $0{,}125$ jako zlomek; najdi číslo mezi $\frac14$ a $\frac13$.
\item Rozhodni, zda $\frac{3+9}{3+15}=\frac9{15}$. Ověř oba zlomky
zkrácením; pojmenuj chybu. Rozhodni také, zda $\frac04$ a $\frac40$ označují čísla.
\end{enumerate}
\tcblower
1. $\frac23$, $-\frac57$, $\frac47$.
2. $-\frac{35}{60}$, $\frac{27}{42}$.
3. $\frac{45}{72}>\frac{44}{72}$; $-\frac{21}{30}<-\frac{20}{30}$.
4. $-\frac58<-\frac12<0{,}7<\frac34$.
5. $\frac18$; například $\frac7{24}$ mezi $\frac6{24}$ a $\frac8{24}$.
6. Ne: vlevo $\frac{12}{18}=\frac23$, vpravo $\frac35$; chybně se krátí
sčítance. $\frac04=0$ je číslo, $\frac40$ není definováno.
\end{zadaniReseniBox}

\end{document}
